\section{Long Context: Kimi Linear y Delta Attention}

\subsection{Por qué el contexto largo es fundamental en la era de los agentes}

Yang Zhilin (杨植麟) citó una gráfica clásica, del nivel de una «joya oculta»: la comparación entre Transformer y LSTM. El Transformer no solo obtiene una loss menor con los mismos parámetros y el mismo número de tokens, sino que, de manera más importante, es capaz de \textbf{mejorar de manera continua su predicción por medio del contexto}: a medida que aumenta el índice de tokens, la loss del Transformer sigue bajando, mientras que la del LSTM se satura después de cierto rango.

\begin{knowledgebox}{La verdadera ventaja del Transformer}
El Transformer no gana solo por ser «más fuerte», sino porque puede aprovechar un contexto más largo para obtener una ganancia de información continua. Esta característica es especialmente crucial en la era de los agentes: un agente necesita ejecutarse durante días o incluso semanas para completar tareas complejas (como escribir un kernel de Linux desde cero), lo cual exige que el modelo procese de manera eficiente una trayectoria extremadamente larga.
\end{knowledgebox}

\subsection{Kimi Delta Attention: decaimiento de grano fino}

La innovación central de la arquitectura Kimi Linear es \textbf{Kimi Delta Attention}, que mejora el mecanismo de memoria recurrente de GDR (Generalized Delta Rule).

La linear attention tradicional usa un factor de decaimiento escalar global (global scalar decay), lo cual obliga a elegir entre «olvidarlo todo» y «recordarlo todo». Kimi Delta Attention introduce un \textbf{factor de decaimiento de grano fino} (fine-grained decay):

$$\alpha \in \mathbb{R}^{d \times d} \quad (\text{matriz diagonal, no un escalar})$$

De este modo, distintos canales pueden tener distintas velocidades de decaimiento:
\begin{itemize}
    \item algunos canales decaen lentamente, con lo que conservan información de largo alcance
    \item algunos canales decaen rápidamente, con lo que olvidan a tiempo y absorben información nueva
\end{itemize}

\subsection{Implementación eficiente: fórmula por chunks e inversión de matrices}

Para paralelizar en las GPU actuales, es necesario convertir la fórmula recurrente en una forma por chunks (chunk-wise). Pero la matriz diagonal $\alpha$ no permite extraer un factor común con la misma facilidad que un escalar, lo cual plantea un enorme desafío de ingeniería.

\begin{importantbox}{Una fórmula paralela exactamente equivalente}
Al introducir una \textbf{operación de inversión de matrices} y un \textbf{factor de decaimiento acumulado}, el equipo reescribió la fórmula recurrente como tres ecuaciones que se pueden calcular en paralelo. No se trata de una aproximación: es una fórmula matemáticamente equivalente en sentido estricto, que logra un paralelismo eficiente sin sacrificar precisión alguna.
\end{importantbox}

\subsection{Comparación de desempeño}

Kimi Linear obtuvo una ventaja integral en comparaciones justas:
\begin{itemize}
    \item \textbf{Tareas de contexto corto} (MMAU): superior a MLA y a GDN
    \item \textbf{Tareas de contexto largo} (Ruler): también superior a las demás variantes, y con mayor eficiencia
    \item Al escalar a un millón de tokens o más, la ventaja en eficiencia se vuelve todavía más notable
    \item Es la \textbf{primera arquitectura que supera a full attention en todas las dimensiones} (contexto corto, entrada larga, salida larga)
\end{itemize}

La arquitectura combina linear attention con full attention en una proporción de 3:1, con lo que logra un equilibrio entre la capacidad de contexto largo y la eficiencia.

\subsection{Resumen de la sección}

Kimi Linear, mediante el factor de decaimiento de grano fino y una fórmula paralela exactamente equivalente, logra un desempeño mejor y una eficiencia mayor que full attention, con una ventaja especialmente notable en escenarios de contexto extremadamente largo.

